% ═══════════════════════════════════════════════════════════════
% LERNBLATT TEMPLATE (Best-Practice-Design v2.0)
% Spanisch - 
% ═══════════════════════════════════════════════════════════════
% DESIGN-PRINZIPIEN:
% - Sans-Serif (Helvetica): Fließtext
% - Serif (Palatino): Spanische Vokabeln (visuelle Distinktion)
% - Farbe NUR didaktisch begründet (Stammwechsel, Farbvokabeln)
% - Minimale Rahmen (kein "visual noise")
% ═══════════════════════════════════════════════════════════════

\documentclass[11pt, a4paper]{article}

\usepackage[utf8]{inputenc}
\usepackage[T1]{fontenc}
\usepackage[ngerman]{babel}

% --- SCHRIFTEN ---
% Helvetica (sans-serif) als Grundschrift
\usepackage{helvet}
\renewcommand{\familydefault}{\sfdefault}
% Palatino (serif) für spanische Vokabeln
\usepackage{palatino}
\newcommand{\seriftext}[1]{{\fontfamily{ppl}\selectfont #1}}

\usepackage{microtype}
\usepackage[margin=1.8cm, top=2cm, bottom=1.5cm]{geometry}
\usepackage{enumitem}
\usepackage{tabularx}
\usepackage{booktabs}
\usepackage{array}
\usepackage{multicol}
\usepackage{tikz}

\usepackage{fancyhdr}
\usepackage{lastpage}

% ═══════════════════════════════════════════════════════════════
% VARIABLEN (ANPASSEN!)
% ═══════════════════════════════════════════════════════════════

\newcommand{\lektion}{X}                    % Lektionsnummer
\newcommand{\thema}{Thema}                  % Spanischer Titel
\newcommand{\untertitel}{Deutscher Untertitel}
\newcommand{\klasse}{8}                     % Klassenstufe
\newcommand{\lernziele}{Lernziel 1 · Lernziel 2 · Lernziel 3}

% ═══════════════════════════════════════════════════════════════
% FARBEN
% ═══════════════════════════════════════════════════════════════

\usepackage{xcolor}

% Didaktische Farben (NUR für Stammwechsel, Farbvokabeln etc.)
\definecolor{colorrojo}{HTML}{C62828}       % Rot: Stammwechsel e→ie, o→ue
\definecolor{colorazul}{HTML}{1565C0}       % Blau
\definecolor{colorverde}{HTML}{2E7D32}      % Grün
\definecolor{coloramarillo}{HTML}{F9A825}   % Gelb
\definecolor{colormarron}{HTML}{6D4C41}     % Braun
\definecolor{colorgris}{HTML}{757575}       % Grau
\definecolor{colornegro}{HTML}{222222}      % Schwarz
\definecolor{colorblanco}{HTML}{FFFFFF}     % Weiß

% Strukturfarben (dezent)
\definecolor{hellgrau}{HTML}{F5F5F5}
\definecolor{rahmen}{HTML}{CCCCCC}
\definecolor{akzent}{HTML}{666666}

% ═══════════════════════════════════════════════════════════════
% BOXEN (minimalistisch)
% ═══════════════════════════════════════════════════════════════

\usepackage{tcolorbox}
\tcbuselibrary{skins}

% Lernziel-Box (hellgrauer Hintergrund, kein Rahmen)
\newtcolorbox{lernzielbox}{
    colback=hellgrau, colframe=hellgrau, boxrule=0pt, arc=0pt,
    left=10pt, right=10pt, top=8pt, bottom=8pt,
    before skip=6pt, after skip=8pt
}

% Grammatik-Box (Randlinie links)
\newtcolorbox{grammatikbox}[1][]{
    colback=white, colframe=white, boxrule=0pt, arc=0pt,
    left=10pt, right=6pt, top=6pt, bottom=6pt,
    borderline west={3pt}{0pt}{akzent},
    before skip=4pt, after skip=4pt
}

% Merk-Box (Linie oben)
\newtcolorbox{merkbox}[1][]{
    colback=hellgrau, colframe=hellgrau, boxrule=0pt, arc=0pt,
    left=10pt, right=10pt, top=8pt, bottom=8pt,
    borderline north={2pt}{0pt}{colornegro},
    before skip=6pt, after skip=6pt
}

% Vokabel-Box (dezenter Rahmen)
\newtcolorbox{vokabelbox}[1][]{
    colback=white, colframe=rahmen, boxrule=0.5pt, arc=0pt,
    left=8pt, right=8pt, top=6pt, bottom=6pt,
    before skip=4pt, after skip=4pt
}

% Tipp-Box (Randlinie links, dezent)
\newtcolorbox{tippbox}{
    colback=white, colframe=white, boxrule=0pt, arc=0pt,
    left=8pt, right=4pt, top=3pt, bottom=3pt,
    borderline west={2pt}{0pt}{akzent}
}

% Zusatz-Box (Linie oben)
\newtcolorbox{zusatzbox}[1][]{
    colback=white, colframe=rahmen, boxrule=0.5pt, arc=0pt,
    left=8pt, right=8pt, top=6pt, bottom=6pt,
    borderline north={1.5pt}{0pt}{akzent}
}

% ═══════════════════════════════════════════════════════════════
% BEFEHLE
% ═══════════════════════════════════════════════════════════════

% Spanisch: Serif + fett (visuelle Distinktion)
\newcommand{\es}[1]{\seriftext{\textbf{#1}}}

% Deutsch: kursiv + Akzentfarbe
\newcommand{\de}[1]{\textit{\textcolor{akzent}{#1}}}

% Grammatikbegriff: Kapitälchen
\newcommand{\gram}[1]{\textsc{#1}}

% AB-Verweis
\newcommand{\abmark}[1]{{\footnotesize\textcolor{akzent}{$\rightarrow$ AB #1}}}

% Abschnittstitel
\newcommand{\abschnitt}[2]{\noindent\textbf{\large #1. #2}}

% Stammwechsel hervorheben (rot)
\newcommand{\stamm}[1]{\textcolor{colorrojo}{#1}}

% Farbmuster für Farbvokabeln
\newcommand{\farbmuster}[1]{%
    \tikz[baseline=-0.5ex]\fill[#1] (0,0) rectangle (8pt,8pt);%
}

\setlength{\parindent}{0pt}
\setlength{\parskip}{5pt}

% ═══════════════════════════════════════════════════════════════
% KOPF- UND FUSSZEILE
% ═══════════════════════════════════════════════════════════════

\pagestyle{fancy}
\fancyhf{}
\renewcommand{\headrulewidth}{0.5pt}
\renewcommand{\footrulewidth}{0pt}
\lhead{\footnotesize\textsc{Spanisch} · Klasse \klasse}
\rhead{\footnotesize\thema}
\cfoot{\footnotesize\thepage\,/\,\pageref{LastPage}}

% ═══════════════════════════════════════════════════════════════
% DOKUMENT
% ═══════════════════════════════════════════════════════════════

\begin{document}

% --- TITEL ---
\begin{center}
    {\LARGE\bfseries \thema}\\[3pt]
    {\small \untertitel}
\end{center}

\vspace{4pt}

% --- EINLEITUNG (motivierend, 2-3 Sätze) ---
% HIER: Motivierender Einleitungstext mit Alltagsbezug
Stell dir vor, du bist in Spanien und...
Der Verkäufer fragt: \es{¿...?} -- \de{...?}

\vspace{2pt}

% --- LERNZIELE ---
\begin{lernzielbox}
\textbf{Das lernst du:}\quad \lernziele
\end{lernzielbox}

% ═══════════════════════════════════════════════════════════════
% ZWEI SPALTEN
% ═══════════════════════════════════════════════════════════════

\begin{multicols}{2}

% ---------------------------------------------------------------
% ABSCHNITT 1
% ---------------------------------------------------------------
\abschnitt{1}{Título español} \de{-- Deutsche Übersetzung}

\vspace{4pt}

\begin{tabular}{@{}ll@{}}
\es{palabra} & \de{Wort} \\[0.4em]
\es{otra palabra} & \de{anderes Wort} \\[0.4em]
% weitere Vokabeln...
\end{tabular}

\vspace{4pt}

\begin{tippbox}
\footnotesize\textit{Tipp:} Merkhilfe oder interessanter Fakt.
\end{tippbox}

\abmark{1}

\vspace{8pt}

% ---------------------------------------------------------------
% ABSCHNITT 2
% ---------------------------------------------------------------
\abschnitt{2}{Título español} \de{-- Deutsche Übersetzung}

\vspace{4pt}

\begin{grammatikbox}
\textbf{\gram{Grammatikregel}}

Erklärung der Regel...\\
{\small Beispiel: \es{...} · \es{...}}
\end{grammatikbox}

\abmark{2}

\columnbreak

% ---------------------------------------------------------------
% ABSCHNITT 3: KONJUGATION
% ---------------------------------------------------------------
\abschnitt{3}{Verbos} \de{-- Verben}

\vspace{4pt}

Mit diesen Verben sagst du, was du \textbf{willst}, \textbf{kannst} oder \textbf{musst}.

\begin{grammatikbox}
{\small
\begin{tabular}{@{}l@{\,}l@{\,}l@{}}
& \textbf{querer} & \textbf{poder} \\
& \de{wollen} & \de{können} \\[0.3em]
\midrule
yo & \es{qu\stamm{ie}ro} & \es{p\stamm{ue}do} \\
tú & \es{qu\stamm{ie}res} & \es{p\stamm{ue}des} \\
él/ella & \es{qu\stamm{ie}re} & \es{p\stamm{ue}de} \\
nosotros & \es{queremos} & \es{podemos} \\
vosotros & \es{queréis} & \es{podéis} \\
ellos & \es{qu\stamm{ie}ren} & \es{p\stamm{ue}den} \\
\end{tabular}}
\end{grammatikbox}

\vspace{2pt}
{\footnotesize\textit{Merke:} Stammwechsel \stamm{e→ie} und \stamm{o→ue} -- außer bei nosotros/vosotros!}

\abmark{3}

\vspace{8pt}

% ---------------------------------------------------------------
% ABSCHNITT 4: DIALOG
% ---------------------------------------------------------------
\abschnitt{4}{Comunicación} \de{-- Dialog}

\vspace{4pt}

\textbf{Vendedor/a:}\\
{\small
\es{¿En qué puedo ayudarte?}\\
\es{¿Qué color quieres?}}

\vspace{4pt}

\textbf{Cliente:}\\
{\small
\es{Quiero comprar...}\\
\es{¿Cuánto cuesta?}}

\abmark{6}

\end{multicols}

% ═══════════════════════════════════════════════════════════════
% SEITENUMBRUCH
% ═══════════════════════════════════════════════════════════════
\newpage

% ---------------------------------------------------------------
% ABSCHNITT 5: HAUPTTHEMA
% ---------------------------------------------------------------
\abschnitt{5}{Tema principal} \de{-- Hauptthema}

\vspace{4pt}

Einleitender Text zum Hauptthema...

\vspace{4pt}

\begin{merkbox}
\textbf{Die wichtigsten Punkte}

\vspace{4pt}
\begin{tabular}{@{}rl@{\hspace{0.6cm}}rl@{\hspace{0.6cm}}rl@{}}
100 & \es{cien} & 400 & \es{cuatrocientos} & 700 & \es{setecientos} \\[0.3em]
200 & \es{doscientos} & 500 & \es{quinientos} & 800 & \es{ochocientos} \\[0.3em]
300 & \es{trescientos} & 600 & \es{seiscientos} & 900 & \es{novecientos} \\
\end{tabular}
\end{merkbox}

\vspace{4pt}

\begin{multicols}{2}

\begin{grammatikbox}
\textbf{Drei Regeln}

\textbf{1.} Erste Regel...\\
\phantom{1.} Erklärung mit \es{Beispiel}

\vspace{3pt}
\textbf{2.} Zweite Regel...\\
\phantom{2.} Erklärung

\vspace{3pt}
\textbf{3.} Dritte Regel...\\
\phantom{3.} Erklärung
\end{grammatikbox}

\columnbreak

\textbf{Beispiele:}\\[2pt]
{\small
\begin{tabular}{@{}rl@{}}
350 & = \es{trescientos cincuenta} \\[0.3em]
725 & = \es{setecientos veinticinco} \\[0.3em]
999 & = \es{novecientos noventa y nueve} \\
\end{tabular}}

\vspace{8pt}

{\footnotesize
\textit{Basis:} Beginne mit den einfachen Beispielen.\\
\textit{Fortgeschritten:} Übe die komplexeren Formen.}

\end{multicols}

\abmark{4 + 5}

\vspace{8pt}

% ═══════════════════════════════════════════════════════════════
% ZUSATZAUFGABEN
% ═══════════════════════════════════════════════════════════════

\begin{zusatzbox}
{\small\textbf{Zusatzaufgaben für Schnelle}}

\vspace{4pt}

\begin{multicols}{2}
\footnotesize

\textbf{Z1: Beschreibung}\\
Beschreibe 5 Elemente mit Adjektiven auf Spanisch.

\vspace{4pt}

\textbf{Z2: Partner-Dialog}\\
Schreibt einen Dialog (8 Zeilen).

\columnbreak

\textbf{Z3: Rechnen}\\
Ergebnis als spanisches Wort:\\
a) 350 + 275 = \rule{2cm}{0.3pt}\\
b) 1000 $-$ 456 = \rule{2cm}{0.3pt}

\vspace{4pt}

\textbf{Z4: Kreativ}\\
Freie Aufgabe zur Vertiefung.

\end{multicols}
\end{zusatzbox}

\vspace{8pt}

% ═══════════════════════════════════════════════════════════════
% KULTURELLES (optional)
% ═══════════════════════════════════════════════════════════════

\begin{grammatikbox}
\textbf{¿Sabías que...?} \de{-- Wusstest du?}

\vspace{4pt}
\begin{multicols}{2}
\footnotesize

$\bullet$ Interessanter kultureller Fakt über Spanien oder Lateinamerika.

\vspace{3pt}
$\bullet$ Sprachlicher Hinweis oder Unterschied zum Deutschen.

\columnbreak

$\bullet$ Alltagsrelevante Information für Spanien-Reisende.

\vspace{3pt}
$\bullet$ Grammatik-Tipp oder Eselsbrücke.

\end{multicols}
\end{grammatikbox}

\end{document}
